\chapter{О чём молчат фигурные скобки?}
\label{ch:first-steps}

\begin{quote}
  \textit{«Математика — это искусство давать одинаковые имена разным вещам».}
  \begin{flushright}
    --- Анри Пуанкаре
  \end{flushright}
\end{quote}

\section{Чайный парадокс, с которого всё началось}
\label{sec:tea-paradox}

Помню, как в детстве я поспорил с другом, что такое «куча». Мы сидели на даче, перебирали камушки. Я сказал: «Куча — это когда камней больше десяти». Он усмехнулся и убрал один камень: «А теперь? Девять — уже не куча?» Я замялся. А если ровно десять? В общем, мы быстро поняли, что слово «куча» какое-то скользкое. С ним нельзя рассуждать строго.

Примерно так же философы и математики много веков спорили о множествах. Что значит «множество всех красивых девушек»? А «множество всех множеств»? В последнем случае вообще случилась катастрофа — парадокс Рассела. Представь себе деревенского парикмахера, который бреет всех тех и только тех жителей деревни, кто не бреется сам. Вопрос: бреет ли он самого себя? Если бреет — то не должен, а если не бреет — то обязан. Взрыв мозга.

Математики поняли: чтобы строить надёжный фундамент, нужно отказаться от расплывчатых «куч» и говорить только о таких собраниях объектов, для которых можно чётко сказать: этот элемент принадлежит множеству, а этот — нет. Не по настроению, а по строгому закону. И записывать этот закон мы будем с помощью волшебных фигурных скобок.

Вот с них мы и начнём. Как превратить интуитивное «много предметов» в точный инструмент?

\section{Говорим на языке свойств}
\label{sec:language-properties}

Вместо того чтобы перечислять всё что угодно, математик говорит: «Возьмём множество всех $x$, для которых выполняется свойство $P(x)$». Запись $\{x \mid P(x)\}$ — это как заклинание. Оно означает: собери все объекты, которые обладают свойством $P$. Например, $\{x \mid x \text{ — натуральное число и } x < 5\}$ — это $\{1,2,3,4\}$. Просто и ясно.

Но есть одно \textbf{железное правило}: свойство $P(x)$ должно быть таким, чтобы для любого объекта мы могли однозначно решить, истинно оно или ложно. Никаких «красивых девушек», если мы не можем записать это условие математически. Иначе снова здравствуй, парадокс Рассела.

Главное отношение в теории множеств — \emph{принадлежность}. Мы пишем $x \in A$, если $x$ является элементом $A$, и $x \notin A$ в противном случае. И всё. Множество задано, если для любого $x$ мы знаем ответ на вопрос $x \in A$?.

% Возможно, здесь уместно небольшое замечание про пустое множество
\smallskip
\textbf{Пустое множество.} Особый случай: $\emptyset = \{\,\}$. В нём нет ни одного элемента. Это не «ничто», а коробка, в которой ничего не лежит. И оно одно-единственное, потому что его полностью определяет свойство $x \neq x$.

\section{Погружение в задачу: Одинаковы ли эти множества?}
\label{sec:same-sets}

Давай сразу нырнём в воду. Пусть у нас есть два списка: $A = \{1, 2, 3\}$ и $B = \{2, 1, 3\}$. Это одно и то же множество или разное? Интуитивно: конечно одно. Пакет с тремя конфетами «Мишка», «Белочка» и «Грильяж» — это тот же пакет, даже если мы переложили конфеты. Порядок не важен. Но как это доказать? Вдруг найдётся дотошный инопланетянин, который скажет: «Нет, первый элемент $1$, а в $B$ первый элемент $2$ — они отличаются!»

Первая мысль: «Ну, мы можем переставить элементы». Но перестановка — это уже действие с элементами, а нам нужно свойство самих множеств. Тупик: если бы мы определяли множество как список, где порядок важен, мы бы получили совсем другую теорию (кортежи). Но мы же договорились, что множество — это только ответ на вопрос «принадлежит или нет».

Значит, чтобы доказать равенство $A$ и $B$, мы должны убедиться: любой элемент из $A$ лежит в $B$, и наоборот. Смотрим: $1 \in A$, проверим $1 \in B$? Да. $2 \in A \to 2 \in B$. $3 \in A \to 3 \in B$. Теперь в обратную сторону: $2 \in B \to 2 \in A$, $1 \in B \to 1 \in A$, $3 \in B \to 3 \in A$. Всё сошлось. Принцип \textbf{экстенсиональности}: множества равны, если состоят из одних и тех же элементов. И точка. Никаких «первых» и «вторых» номеров.

\begin{axiom}[Экстенсиональность]
  $A = B$ тогда и только тогда, когда $\forall x \,(x \in A \iff x \in B)$.
\end{axiom}

\section{Момент озарения: круги на песке}
\label{sec:venn-insight}

На этом этапе мы обычно рисуем круги на доске и закрашиваем пересечения. Диаграммы Эйлера--Венна — гениальное изобретение. Они помогают нам увидеть операции наглядно. Пересечение $A \cap B$ — это общая часть кругов. Объединение $A \cup B$ — всё, что закрашено хоть как-то. Разность $A \setminus B$ — из $A$ убрали кусочек $B$.

Но попробуй нарисовать пустое множество $\emptyset$. Просто ничего. А потом доказать с помощью кругов, что $(A \cup B) \cap C = (A \cap C) \cup (B \cap C)$. Ты утонишь в штриховках и проверке областей. Интуиция кругов прекрасна, но для строгих доказательств нам понадобится язык логики. И вот тут происходит озарение: каждая операция — это просто логическая связка.

\begin{itemize}
  \item $x \in A \cap B$ \quad означает \quad $x \in A$ \textbf{и} $x \in B$.
  \item $x \in A \cup B$ \quad означает \quad $x \in A$ \textbf{или} $x \in B$.
  \item $x \in A \setminus B$ \quad означает \quad $x \in A$ \textbf{и не} $x \in B$.
  \item $x \in A \triangle B$ \quad означает \quad ($x \in A$ \textbf{и не} $x \in B$) \textbf{или} ($x \in B$ \textbf{и не} $x \in A$).
\end{itemize}

А дальше мы можем просто жонглировать этими И, ИЛИ, НЕ, используя свойства логики, которые мы знаем. Это как вместо того, чтобы переставлять мебель в комнате, написать план комнаты и двигать карандашные линии.

\section{Формализация: азбука операций}
\label{sec:operations-formal}

Давай запишем строго, но без занудства. Пусть $U$ — некоторый универсум (мир, где лежат все интересующие нас объекты). Тогда для любых подмножеств $A, B \subseteq U$:

\begin{align}
  A \cup B &= \{x \in U \mid x \in A \lor x \in B\},\\
  A \cap B &= \{x \in U \mid x \in A \land x \in B\},\\
  A \setminus B &= \{x \in U \mid x \in A \land \lnot (x \in B)\},\\
  A \triangle B &= (A \setminus B) \cup (B \setminus A).
\end{align}

Теперь докажем первый попавшийся закон — дистрибутивность:
\[
A \cup (B \cap C) = (A \cup B) \cap (A \cup C).
\]
Я мог бы нарисовать три круга и старательно раскрашивать. А могу рассуждать так: возьмём произвольный $x$. Что значит, что $x$ принадлежит левой части? Это значит: $x \in A \lor (x \in B \land x \in C)$. По свойствам логики (дистрибутивность $\lor$ относительно $\land$) это эквивалентно: $(x \in A \lor x \in B) \land (x \in A \lor x \in C)$. А это в точности запись $x \in (A \cup B) \cap (A \cup C)$. Готово. Никаких картинок, только цепочка равносильностей.

\subsection{Как доказывать равенство множеств: урок дистрибутивности}
\label{sec:distrib-proof}

Теперь разберёмся с задачей, которую ты прислал:
\[
(A \cap B) \cup C = (A \cup C) \cap (B \cup C).
\]
Формула красивая, почти заклинание. Но как подступиться?

\textbf{Наивный порыв: «Я нарисую круги!»} Первое, что приходит в голову — нарисовать три пересекающихся круга на салфетке и заштриховать. Я так и сделал. Картинки совпали! Можно выдохнуть? Нет. Картинка — только подсказка, а не доказательство. Мы хотим научиться доказывать строго.

\textbf{Тупик: «Преобразую как в алгебре»}. Иногда начинаешь чисто механически переписывать: «А∩В∪С… порядок действий нарушен, скобки путаются…» Бесполезно.

\textbf{Первый шаг к прозрению:} что вообще значит «равенство множеств»? Чтобы доказать $M = N$, нужно установить два включения:
\begin{enumerate}
  \item $M \subseteq N$ (все элементы $M$ лежат в $N$),
  \item $N \subseteq M$ (все элементы $N$ лежат в $M$).
\end{enumerate}
Значит, наша задача распадается на две половинки. Уже легче.

\textbf{Часть 1:} $(A \cap B) \cup C \subseteq (A \cup C) \cap (B \cup C)$.
Берём произвольный $x \in (A \cap B) \cup C$. По определению объединения это значит: $x \in A \cap B$ \textbf{или} $x \in C$. Перед нами развилка.
\begin{itemize}
  \item \textit{Случай $x \in A \cap B$.} Тогда $x \in A$ и $x \in B$. Но если $x \in A$, то $x \in A \cup C$. Аналогично $x \in B \cup C$. Значит, $x$ лежит в пересечении $(A \cup C) \cap (B \cup C)$.
  \item \textit{Случай $x \in C$.} Тогда $x \in A \cup C$ (потому что всё, что в $C$, уже входит в объединение) и $x \in B \cup C$. Снова $x$ в пересечении.
\end{itemize}
Оба случая дают нужный результат. Включение доказано.

\textbf{Часть 2:} $(A \cup C) \cap (B \cup C) \subseteq (A \cap B) \cup C$.
Берём произвольный $x \in (A \cup C) \cap (B \cup C)$. Значит, $x \in A \cup C$ и $x \in B \cup C$. Теперь ключевой вопрос: находится ли $x$ в $C$? Ответов только два: да или нет.
\begin{itemize}
  \item \textit{Альтернатива 1: $x \in C$.} Тогда немедленно $x \in (A \cap B) \cup C$ (достаточно того, что он в $C$).
  \item \textit{Альтернатива 2: $x \notin C$.} Раз $x \in A \cup C$, но $x \notin C$, то обязательно $x \in A$. Из $x \in B \cup C$ и $x \notin C$ получаем $x \in B$. Таким образом, $x \in A \cap B$, а значит и в $(A \cap B) \cup C$.
\end{itemize}
Снова оба пути ведут в цель. Включение доказано.

Два доказанных включения по определению означают равенство множеств. $\square$

Видишь? Доказательство как детектив: мы не просто предъявляем цепочку символов, а разыгрываем сценарий «Что если элемент здесь? А что если там?». Главный приём — \textbf{расслоение на случаи}. Именно он превращает туманную картинку в прозрачную логику.

\section{Ряд упражнений}
\label{sec:exercises-ch1}

Теперь твоя очередь. Попробуй доказать следующие тождества, используя только определения и логические рассуждения (без кругов!). Ниже даны не решения, а \textbf{подсказки-размышления}.

\begin{enumerate}
  \item \emph{Закон де Моргана:} $U \setminus (A \cup B) = (U \setminus A) \cap (U \setminus B)$. \\
    \textit{Подсказка.} «Не (в $A$ или в $B$)» звучит двусмысленно. А что если перевести: «Элемент не принадлежит объединению» означает, что он не лежит ни в $A$, ни в $B$. Как это записать через «и» и «не»? Попробуй пойти от логики высказываний.

  \item \emph{Ассоциативность пересечения:} $A \cap (B \cap C) = (A \cap B) \cap C$. \\
    \textit{Подсказка.} Возьми элемент $x$, который лежит в левой части. Какие три утверждения о его принадлежности ты можешь записать? А теперь перегруппируй скобки в логическом «и» — оно ведь ассоциативно.

  \item \emph{Ещё одна дистрибутивность:} $A \cap (B \triangle C) = (A \cap B) \triangle (A \cap C)$. \\
    \textit{Подсказка.} Симметрическая разность — это «исключающее или». Придётся разбирать случаи, но можно и иначе: переведи всё на язык характеристических функций (мы поговорим о них в следующей главе, но вдруг ты уже знаком). А пока честно рассмотри, что значит $x \in A \cap (B \triangle C)$. Это когда $x \in A$ и $x$ принадлежит ровно одному из $B$, $C$. Убедись, что ровно то же получается и в правой части.

  \item \emph{Пустое множество:} Докажи, что $\emptyset \subseteq A$ для любого $A$. \\
    \textit{Подсказка.} Определение подмножества: $X \subseteq Y$, если $\forall x \,(x \in X \Rightarrow x \in Y)$. Что делает импликация с ложной посылкой? У пустого множества нет элементов, значит посылка $x \in \emptyset$ всегда ложна…

  \item \emph{Равенство через включения:} Пусть $A = \{1,2,3\}$, $B = \{x \in \mathbb{N} \mid x \le 3\}$, $C = \{3,1,2\}$. Покажи, что $A = B = C$, используя только принцип экстенсиональности. \\
    \textit{Подсказка.} Выпиши элементы каждого множества и убедись, что они одинаковы. Обрати внимание: порядок записи не важен!
\end{enumerate}

\section{Итоговый вывод}
\label{sec:conclusion-ch1}

Итак, мы научились говорить на точном языке свойств. Множество — это не мешок с картошкой, а чётко заданное условие. Операции объединения, пересечения и разности — это просто перевод логических связок «и», «или», «не» на язык фигурных скобок. Главный инструмент доказательства равенств — включения и разбор случаев.

Это была наша первая формальная глава. Мы подготовили почву для того, чтобы задаться по-настоящему глубокими вопросами: а сколько бывает элементов у множества? Как сравнивать бесконечные совокупности? И что будет, если посмотреть на множество как на функцию, выдающую $0$ или $1$? Этим мы займёмся в следующей главе, где в игру вступит характеристическая функция — наш скромный герой, который превратит операции над множествами в простую арифметику.